% ── Interlude 02: The Instrument of Transformation ───────────
% Object class: tool that converts signal into structure
% Appears after Part II (Mackandal) before Part III (Payne)
% No named characters. No date. No location.

\begin{interlude}{The Instrument of Transformation}

Something enters. Something different leaves.

This is the basic description of any instrument, and it is
insufficient. It says nothing about what is preserved in the
passage, what is lost, what is added by the instrument itself.
Every transformation has a signature---a characteristic distortion,
a preferred direction, a set of inputs it handles well and a set
it does not handle at all. To use an instrument without knowing
its signature is to mistake the instrument's contribution for a
property of what was transformed.

Most instruments are used this way.

The transformation is trusted because the instrument is familiar.
What comes out is taken as a faithful rendering of what went in,
modified only in the ways that were intended. The instrument is
treated as transparent---a pure medium through which the thing
passes unchanged except in form.

It is never transparent.

\medskip

There are instruments that transform by subtraction.

They remove what is not needed, or what is not permitted, or what
cannot survive the passage. What emerges is cleaner than what
entered---less ambiguous, more manageable, easier to transmit. The
cost is not always visible. It appears later, when something is
needed that the instrument discarded, when the question asked of
the output cannot be answered because the relevant information did
not survive the transformation.

There are instruments that transform by addition.

They impose structure on what passes through them---a frame, a
scale, a set of categories that were not present in the input.
What emerges carries the instrument's assumptions embedded within
it, indistinguishable from what was there before. Such instruments
are difficult to identify, because their contribution looks like
discovery rather than imposition.

There are instruments that transform by relation.

They do not act on the thing directly but place it in contact with
something else---a standard, a field, a known quantity---and
measure what results. The output is not the thing itself but its
behaviour under that contact. This is perhaps the most honest form
of transformation: it does not pretend to reveal the thing as it
is, only as it responds.

\medskip

The instrument changes with use.

Not always visibly, not always quickly, but the passage of
material through it leaves traces. Residue accumulates. Tolerances
shift. What the instrument does to its ten-thousandth input is not
quite what it did to its first, even if the inputs are identical.

This is rarely accounted for.

The assumption is that the instrument remains constant while the
inputs vary---that it is the stable term in the relation, the
fixed point against which change is measured. But the instrument
is also a material thing, subject to the same pressures it applies
to everything else.

An instrument that has transformed many things carries within it
the history of those transformations. It has been shaped by what
has passed through it. Its current behaviour is partly a function
of its original design and partly a function of everything it has
already done.

To read the output correctly, one would need to know the
instrument's history.

This knowledge is almost never available.

\medskip

Some transformations cannot be reversed.

What enters in one form does not return to that form when the
process is run backward. The direction matters. The instrument
knows, in some sense, which way time is moving---which end is the
input and which is the output---and this asymmetry is not accidental.
It is a property of what transformation means.

A thing that has been transformed is not the same thing it was.

This seems obvious, but its implications are not always followed.
The transformed thing is treated as equivalent to the original,
just in a different form. The form is treated as superficial,
something that can be changed again without loss. But each
transformation selects, distorts, embeds assumptions, and leaves
the instrument's signature in the result.

To transform something is to commit it to a particular history.

The instrument makes this commitment on behalf of whatever passes
through it, without asking, without recording what was given up.

It only produces the output.

It does not explain what the output cost.

\end{interlude}
